% Options for packages loaded elsewhere
% Options for packages loaded elsewhere
\PassOptionsToPackage{unicode}{hyperref}
\PassOptionsToPackage{hyphens}{url}
\PassOptionsToPackage{dvipsnames,svgnames,x11names}{xcolor}
%
\documentclass[
  letterpaper,
  DIV=11,
  numbers=noendperiod]{scrartcl}
\usepackage{xcolor}
\usepackage{amsmath,amssymb}
\setcounter{secnumdepth}{-\maxdimen} % remove section numbering
\usepackage{iftex}
\ifPDFTeX
  \usepackage[T1]{fontenc}
  \usepackage[utf8]{inputenc}
  \usepackage{textcomp} % provide euro and other symbols
\else % if luatex or xetex
  \usepackage{unicode-math} % this also loads fontspec
  \defaultfontfeatures{Scale=MatchLowercase}
  \defaultfontfeatures[\rmfamily]{Ligatures=TeX,Scale=1}
\fi
\usepackage{lmodern}
\ifPDFTeX\else
  % xetex/luatex font selection
\fi
% Use upquote if available, for straight quotes in verbatim environments
\IfFileExists{upquote.sty}{\usepackage{upquote}}{}
\IfFileExists{microtype.sty}{% use microtype if available
  \usepackage[]{microtype}
  \UseMicrotypeSet[protrusion]{basicmath} % disable protrusion for tt fonts
}{}
\makeatletter
\@ifundefined{KOMAClassName}{% if non-KOMA class
  \IfFileExists{parskip.sty}{%
    \usepackage{parskip}
  }{% else
    \setlength{\parindent}{0pt}
    \setlength{\parskip}{6pt plus 2pt minus 1pt}}
}{% if KOMA class
  \KOMAoptions{parskip=half}}
\makeatother
% Make \paragraph and \subparagraph free-standing
\makeatletter
\ifx\paragraph\undefined\else
  \let\oldparagraph\paragraph
  \renewcommand{\paragraph}{
    \@ifstar
      \xxxParagraphStar
      \xxxParagraphNoStar
  }
  \newcommand{\xxxParagraphStar}[1]{\oldparagraph*{#1}\mbox{}}
  \newcommand{\xxxParagraphNoStar}[1]{\oldparagraph{#1}\mbox{}}
\fi
\ifx\subparagraph\undefined\else
  \let\oldsubparagraph\subparagraph
  \renewcommand{\subparagraph}{
    \@ifstar
      \xxxSubParagraphStar
      \xxxSubParagraphNoStar
  }
  \newcommand{\xxxSubParagraphStar}[1]{\oldsubparagraph*{#1}\mbox{}}
  \newcommand{\xxxSubParagraphNoStar}[1]{\oldsubparagraph{#1}\mbox{}}
\fi
\makeatother


\usepackage{longtable,booktabs,array}
\usepackage{calc} % for calculating minipage widths
% Correct order of tables after \paragraph or \subparagraph
\usepackage{etoolbox}
\makeatletter
\patchcmd\longtable{\par}{\if@noskipsec\mbox{}\fi\par}{}{}
\makeatother
% Allow footnotes in longtable head/foot
\IfFileExists{footnotehyper.sty}{\usepackage{footnotehyper}}{\usepackage{footnote}}
\makesavenoteenv{longtable}
\usepackage{graphicx}
\makeatletter
\newsavebox\pandoc@box
\newcommand*\pandocbounded[1]{% scales image to fit in text height/width
  \sbox\pandoc@box{#1}%
  \Gscale@div\@tempa{\textheight}{\dimexpr\ht\pandoc@box+\dp\pandoc@box\relax}%
  \Gscale@div\@tempb{\linewidth}{\wd\pandoc@box}%
  \ifdim\@tempb\p@<\@tempa\p@\let\@tempa\@tempb\fi% select the smaller of both
  \ifdim\@tempa\p@<\p@\scalebox{\@tempa}{\usebox\pandoc@box}%
  \else\usebox{\pandoc@box}%
  \fi%
}
% Set default figure placement to htbp
\def\fps@figure{htbp}
\makeatother





\setlength{\emergencystretch}{3em} % prevent overfull lines

\providecommand{\tightlist}{%
  \setlength{\itemsep}{0pt}\setlength{\parskip}{0pt}}



 


% === Fuentes ===
\usepackage{fontspec}
\setmainfont{Arial}

% === Interlineado ===
\usepackage{setspace}
\onehalfspacing

% === Colores ===
\usepackage{xcolor}
\definecolor{Pantone253C}{HTML}{AD25A8}

% === Encabezados y pies (KOMA-Script) ===
\usepackage{scrlayer-scrpage}
\clearpairofpagestyles
\pagestyle{scrheadings}
\setkomafont{pagehead}{\sffamily\footnotesize}
\ihead{\textnormal{\color{Pantone253C}\articuloTitulo}}
\ohead{\pagemark}
\setkomafont{pagefoot}{\sffamily\footnotesize}
\ifoot{\textnormal{\color{Pantone253C}An@lítica}}
\rofoot{\textnormal{\color{Pantone253C}Nueva Época ∫∫ Volumen 1, Número 1}} 


\KOMAoptions{headsepline=0.1pt, footsepline=0.1pt}
\addtokomafont{headsepline}{\color{Pantone253C}}
\addtokomafont{footsepline}{\color{Pantone253C}}
\addtokomafont{pagenumber}{\color{Pantone253C}}
\KOMAoption{captions}{tableheading}
\makeatletter
\@ifpackageloaded{caption}{}{\usepackage{caption}}
\AtBeginDocument{%
\ifdefined\contentsname
  \renewcommand*\contentsname{Table of contents}
\else
  \newcommand\contentsname{Table of contents}
\fi
\ifdefined\listfigurename
  \renewcommand*\listfigurename{List of Figures}
\else
  \newcommand\listfigurename{List of Figures}
\fi
\ifdefined\listtablename
  \renewcommand*\listtablename{List of Tables}
\else
  \newcommand\listtablename{List of Tables}
\fi
\ifdefined\figurename
  \renewcommand*\figurename{Figure}
\else
  \newcommand\figurename{Figure}
\fi
\ifdefined\tablename
  \renewcommand*\tablename{Table}
\else
  \newcommand\tablename{Table}
\fi
}
\@ifpackageloaded{float}{}{\usepackage{float}}
\floatstyle{ruled}
\@ifundefined{c@chapter}{\newfloat{codelisting}{h}{lop}}{\newfloat{codelisting}{h}{lop}[chapter]}
\floatname{codelisting}{Listing}
\newcommand*\listoflistings{\listof{codelisting}{List of Listings}}
\makeatother
\makeatletter
\makeatother
\makeatletter
\@ifpackageloaded{caption}{}{\usepackage{caption}}
\@ifpackageloaded{subcaption}{}{\usepackage{subcaption}}
\makeatother
\usepackage{bookmark}
\IfFileExists{xurl.sty}{\usepackage{xurl}}{} % add URL line breaks if available
\urlstyle{same}
\hypersetup{
  pdftitle={Introducción del Editor},
  colorlinks=true,
  linkcolor={blue},
  filecolor={Maroon},
  citecolor={Blue},
  urlcolor={Blue},
  pdfcreator={LaTeX via pandoc}}


\title{Introducción del Editor}
\author{}
\date{}
\begin{document}
\maketitle


\newcommand{\articuloTitulo}{Introducción del Editor}

\section{Experimentación formal en internet o de cómo las Ciencias
Sociales habitan la galaxia: convergencias, algoritmos y narrativas
emergentes}\label{experimentaciuxf3n-formal-en-internet-o-de-cuxf3mo-las-ciencias-sociales-habitan-la-galaxia-convergencias-algoritmos-y-narrativas-emergentes}

La relación entre las Ciencias Sociales y la galaxia internet ha dejado
de ser un objeto de estudio para convertirse en un espacio de
reconfiguración epistemológica, metodológica y política. Los cinco
artículos de investigación aquí reunidos: Capitalismo cognitivo,
gubernamentalidad algorítmica y abordajes de la comunicación de Miguel
Ángel Mata Salazar; Hegemonía algorítmica y visibilidad académica. Una
propuesta metodológica de Christian Hernández Pérez; Nuevas formas de
narrar la realidad. El proceso de individuación y su impacto en la
escritura académica de Paola Vázquez Almanza; Ensayo sobre un pódcast de
investigación sociológica de Fernando Beltrán Nieves; y El podcasting
como herramienta tecnopedagógica en la construcción de una narrativa
transmedia: el caso Versátil de Diego Valdez Serrano, dialogan entre sí
desde distintas aristas de un mismo fenómeno: la transformación radical
de la producción, circulación y legitimación del conocimiento en las
revoluciones permanentes que sufre internet.

Capitalismo cognitivo y algoritmos: el nuevo régimen de visibilidad. El
primer trabajo de este número temático, Capitalismo cognitivo,
gubernamentalidad algorítmica y abordajes de la comunicación, establece
un marco teórico indispensable al vincular el capitalismo cognitivo con
la gubernamentalidad algorítmica, mostrando cómo las lógicas de
acumulación y control se rearticulan en la digitalidad. Esta perspectiva
encuentra un eco empírico en el segundo artículo, Hegemonía algorítmica
y visibilidad académica, que analiza cómo los algoritmos de Google
condicionan la visibilidad de los académicos en México.

Mientras el primero teoriza sobre el poder estructurante de las
plataformas, el segundo lo operacionaliza, revelando un conflicto
central: en un mundo donde la relevancia académica se mide en clics y
posicionamiento, ¿qué sucede con las Ciencias Sociales y Humanidades
críticas, tradicionalmente alejadas de las métricas virales? Ambos
textos convergen en señalar que la academia no es ajena a la
mercantilización de la atención ni a la dataficación de sus prácticas.
Narrativas en transición: de la escritura tradicional al giro sonoro. El
tercer artículo de este número, Nuevas formas de narrar la realidad,
introduce un giro crucial al cuestionar las formas canónicas de
escritura académica, abogando por una mayor imbricación entre
experiencia individual y divulgación. Esta preocupación por democratizar
el conocimiento se materializa en los dos trabajos restantes, dedicados,
ambos, al pódcast. El Ensayo sobre un pódcast de investigación
sociológica analiza las posibilidades analíticas de la sociología en un
medio sonoro, mientras que el segundo artículo, El podcasting como
herramienta tecnopedagógica en la construcción de una narrativa
transmedia: el caso Versátil, explora el formato como herramienta
didáctica en expansión en los recintos universitarios. Los tres textos
comparten una hipótesis implícita: la rigidez de los formatos académicos
tradicionales limita su impacto social, mientras que las narrativas
transmedia ---especialmente el audio--- permiten una circulación ágil y
una apropiación colectiva del saber. Por lo demás, el Ensayo sobre un
pódcast de investigación sociológica explora la técnica de la
interrupción como recurso narrativo, proponiendo un ejercicio
contrafactual sobre Karl Marx que desafía los límites entre ficción y
análisis social. Por su parte, el estudio del pódcast Versátil muestra
cómo este formato se ha convertido en un puente entre la universidad y
el espacio público, siempre que se superen barreras institucionales.
Aquí, la compatibilidad del número temático es evidente: si el
capitalismo cognitivo convierte el conocimiento en commodity, como lo
señala el trabajo de Miguel Ángel Mata, estas iniciativas académicas en
el audio, por el contrario, buscan subvertir esa lógica mediante
prácticas colaborativas y accesibles. Tensiones y desafíos: ¿hacia una
ciencia social expandida?

La relevancia de este número temático radica en su capacidad para mapear
las tensiones contemporáneas de las Ciencias Sociales en la galaxia
internet: desde la dependencia de infraestructuras algorítmicas privadas
hasta la experimentación con formatos que desafían el paper como único
vehículo legítimo. Los cinco textos, en conjunto, trazan un diagnóstico
y una ruta de acción. En efecto, efectúan colectivamente una crítica
sustantiva al denunciar la subsunción de lo académico bajo
racionalidades tecnoliberales (algoritmos, métricas, plataformas), pero
también apuestan por narrativas emergentes (podcasts, técnicas de
interrupción, transmedia) que reconfiguren la relación entre
investigadores y audiencias. Una crítica frontal al explorar y proponer
una suerte de innovación táctica.

Quedan preguntas abiertas. ¿Cómo evitar que la ``visibilidad'' en
internet reproduzca desigualdades ya existentes o, peor aún, las
profundice? ¿Pueden los formatos sonoros eludir la mercantilización que
critica el primer artículo? Lo cierto es que, como lo muestran estos
artículos de investigación, las Ciencias Sociales están obligadas a
navegar ---y transformar--- el ecosistema digital, so pena de volverse
irrelevantes. Esta colección no sólo aporta claves para ese viaje, sino
que lo ejemplifica, plasmando rigor teórico, metodologías audaces y un
compromiso con la circulación plural del conocimiento. En última
instancia, estos textos son artefactos de una misma constelación: la que
busca repensar la academia desde sus fisuras, allí donde internet
desordena jerarquías y obliga a reinventar lo posible. Coronamos este
número temático, por último, con una reseña crítica de Sandra Isabella
Carrillo Cima sobre Todos los museos son novelas de ciencia ficción del
crítico catalán Jorge Carrión. Un trabajo de lectura atenta por parte de
Carrillo Cima sobre las escrituras narrativas en atención a la expansión
inexorable de la inteligencia artificial, una más de las revoluciones
que sufre la galaxia internet.

\begin{flushright}
\textbf{Fernando Beltrán Nieves}\\

\textit{Coordinador Temático}
\end{flushright}




\end{document}
